\subsection{ENGEL'S THEOREM}

Lemma 1. Let \(\mathbf{V}\) be a vector space over \(\mathbf{K}\) . If \(x\) is a nilpotent endomorphism of \(\mathbf{V}\) , the mapping \(y \mapsto [x, y]\) of \(\mathcal{L}(\mathbf{V})\) into \(\mathcal{L}(\mathbf{V})\) is nilpotent.

If \(f\) denotes this mapping, \(f^m(y)\) is a sum of terms of the form \(\pm x^i y x^j\) with \(i + j = m\) . If \(x^k = 0\) , then \(f^{2k - 1}(y) = 0\) for all \(y\) .

THEOREM 1 (Engel). Let V be a vector space over K and g a finite-dimensional subalgebra of \(\mathfrak{gl}(V)\) whose elements are nilpotent endomorphisms of V. If \(V \neq \{0\}\) , there exists \(u \neq 0\) in V such that x.u = 0 for all \(x \in g\) .

The proof proceeds by induction on the dimension \(n\) of \(\mathfrak{g}\) . The theorem is obvious if \(n = 0\) . Suppose that it is true for algebras of dimension \(< n\) .

Let \(\mathfrak{h}\) be a Lie subalgebra of \(\mathfrak{g}\) of dimension \(m < n\) . If \(x \in \mathfrak{h}\) , and \(\operatorname{ad}_{\mathfrak{g}} x\) maps \(\mathfrak{h}\) into itself and defines on passing to the quotient an endomorphism \(\sigma(x)\) of the space \(\mathfrak{g} / \mathfrak{h}\) . By Lemma 1 \(\operatorname{ad}_{\mathfrak{g}} x\) is nilpotent and hence \(\sigma(x)\) is nilpotent. By the induction hypothesis there exists a non-zero element of \(\mathfrak{g} / \mathfrak{h}\) which is annihilated by all the \(\sigma(x)\) , \(x \in \mathfrak{h}\) . It follows that \(\mathfrak{h}\) is an ideal in a certain \((m + 1)\) -dimensional subalgebra of \(\mathfrak{g}\) .

We conclude (by iteration starting with \(\mathfrak{b} = \{0\}\) ) that \(\mathfrak{g}\) has an ideal \(\mathfrak{b}\) of dimension \(n - 1\) . Let \(a \in \mathfrak{g}\) , \(a \notin \mathfrak{b}\) . We again use the induction hypothesis: the \(u \in V\) such that \(x.u = 0\) for all \(x \in \mathfrak{b}\) form a non-zero vector subspace U of V. This subspace is stable under \(a\) ( \(\S 3\) , no. 5, Proposition 5). Since \(a\) is a nilpotent endomorphism of V, there exists a non-zero element of U which is annihilated by \(a\) and hence by every element of \(\mathfrak{g}\) .

COROLLARY 1. For a Lie algebra g to be nilpotent, it is necessary and sufficient that, for all \(x \in \mathfrak{g}\) , ad \(x\) be nilpotent.

The condition is necessary (Proposition 1). Suppose that its sufficiency has been proved for Lie algebras of dimension \(< n\) ( \(n \neq 0\) ). Let \(\mathfrak{g}\) be an \(n\) -dimensional Lie algebra such that, for all \(x \in \mathfrak{g}\) , ad \(x\) is nilpotent. Theorem 1, applied to the set of ad \(x\) ( \(x \in \mathfrak{g}\) ), proves that the centre \(c\) of \(\mathfrak{g}\) is non-zero. Then \(\mathfrak{g}\) is a central extension of the Lie algebra \(\mathfrak{g}/c\) , which is nilpotent by our induction hypothesis. The proof is completed by applying Proposition 2.

COROLLARY 2. Let g be a Lie algebra and h an ideal of g. Suppose that g/h is nilpotent and that, for all \(x \in g\) , the restriction of ad x to h is nilpotent. Then g is nilpotent.

Let \(x \in \mathfrak{g}\) . As \(\mathfrak{g} / \mathfrak{h}\) is nilpotent, there exists an integer \(k\) such that \((\operatorname{ad} x)^k(\mathfrak{g}) \subset \mathfrak{h}\) . By hypothesis there exists an integer \(k'\) such that \((\operatorname{ad} x)^{k'}(\mathfrak{h}) = \{0\}\) . Hence \((\operatorname{ad} x)^{k + k'}(\mathfrak{g}) = \{0\}\) . Corollary 2 is hence a consequence of Corollary 1.

COROLLARY 3. Let V be a vector space and g a finite-dimensional subalgebra of gl(V) whose elements are nilpotent endomorphisms of V. Then g is a nilpotent Lie algebra.

This follows immediately from Lemma 1 and Corollary 1.

Example. The algebra \(\mathfrak{n}(n,\mathbf{K})\) (§ 1, no. 2, Example 2 (3)) is nilpotent.

